\section{Labels and gaps in the circular unitary model}
\label{sec:cue-labels}
We now compute a marked gap statistic in a finite-dimensional model whose
exterior phase is defined explicitly. The argument proves the joint limit
of that phase and the gap; a marginal equidistribution statement would
not suffice for the sign correlation.

Let $U_N$ be Haar distributed in $U(N)$, with angles
$0\le\theta_1<\cdots<\theta_N<2\pi$. Choose an index $I$ uniformly
and independently, after choosing the matrix. For fixed $r\in[0,1)$
and $a\in\mathbb R$, define
\[
 \phi_{N,r}(x)=\Im\sum_{l=1}^N\log(1-r e^{i(x-\theta_l)}),\qquad
 W_{N,i}=i-\frac{N\theta_i}{2\pi}-a-
                   \frac{\phi_{N,r}(\theta_i)}\pi,
 \quad \eta_{N,i}=\operatorname{sgn}\sin(\pi W_{N,i}).
\]
The logarithm is defined by its power series in $r$. Unwrap indices by
$\theta_{i+N}=\theta_i+2\pi$, and set
$D_{j,N}=N(\theta_{I+j}-\theta_I)/(2\pi)$.
Write $k(u)=\sin(\pi u)/(\pi u)$ and
\[
 K^0(x,y)=k(x-y)-k(x)k(y).
\]
This is the kernel of the sine process conditioned to have a point at
zero, with that point removed; this is its reduced Palm process.

\begin{theorem}\label{thm:cue-labels}
For every fixed positive integer $j$, the pair
$(W_{N,I}\bmod2,D_{j,N})$ converges in distribution to $(V,D_j)$,
where $V$ is uniform modulo two and independent of the $j$th positive
point $D_j$ of the reduced Palm sine process. If
$d_2(v)=\min_{m\in\mathbb Z}|v-2m|$, then
\[
 \lim_{N\to\infty}\mathbb P(\eta_{N,I}=1)=\tfrac12,
 \qquad
 \lim_{N\to\infty}\mathbb E(\eta_{N,I}\eta_{N,I+j})
       =\mathbb E\{1-2d_2(j-D_j)\}.
\]
These limits hold uniformly when $r$ ranges in a fixed compact subset
of $[0,1)$.
\end{theorem}

\begin{proof}
Use circle measure $d\theta/(2\pi)$ and let $P_N$ project onto
$1,e^{i\theta},\ldots,e^{i(N-1)\theta}$. Expanding both determinants
in the squared Vandermonde density and integrating coordinate by
coordinate proves
\begin{equation}\label{eq:cue-multiplier}
 \mathbb E\prod_l u(\theta_l)
       =\det(P_NuP_N|_{\operatorname{ran}P_N}).
\end{equation}
Conditioning the marked point to be $x$ replaces $P_N$ by
$Q_{N,x}=P_N-|v_x\rangle\langle v_x|$, where
$v_x(\theta)=P_N(\theta,x)/\sqrt N$. This follows from the Schur
complement of the marked row and column of the correlation determinant.
The marked point is uniform on the circle.

For a finite-rank orthogonal projection $Q$ and $|u|=1$, the positive
contraction $Qu^*(1-Q)uQ$ gives
\begin{equation}\label{eq:cue-leakage}
 |\det(QuQ)|^2
 \le\exp\{-\|(1-Q)uQ\|_{\rm HS}^2\}.
\end{equation}
For $P_N$, counting the Fourier basis vectors that leave its range
under multiplication by a Fourier mode gives
\[
 F_N(u):=\|(1-P_N)uP_N\|_{\rm HS}^2
       =\sum_{n\in\mathbb Z}\min(N,|n|)|\widehat u(n)|^2.
\]
A rank-one removal changes this quantity by at most one: for $P=Q+R$,
$\operatorname{rank}R=1$, the difference is
$-\|RuQ\|_{\rm HS}^2+\|(1-P)uR\|_{\rm HS}^2\in[-1,1]$.

Consider a nonconstant Fourier mode $e^{i\pi kW_{N,I}}$. If $k$ is
even, simultaneous rotation of the matrix leaves the exterior phase
and cyclic gaps invariant and averages $e^{-ikN\theta_I/2}$ to zero.
This conclusion continues to hold after multiplying by any bounded
function of the cyclic gaps. If $k$ is odd, condition on $x$.
Since $I=1+\#\{\theta_l\in(0,x):l\ne I\}$, the mode is, up to a
scalar of modulus one, the determinant in \eqref{eq:cue-multiplier}
with projection $Q_{N,x}$ and multiplier
\[
 u_x(\theta)=(-1)^{\mathbf1_{(0,x)}(\theta)}
     \exp\{-ik\Im\log(1-r e^{i(x-\theta)})\}.
\]
It has two jumps of absolute size two. Two integrations by parts on
its smooth pieces show
\[
 \widehat u_x(n)=\frac{J_0+J_xe^{-inx}}{2\pi in}+O_{x,k,r}(n^{-2}),
 \qquad |J_0|=|J_x|=2.
\]
For fixed $0<x<2\pi$, the cross sum
$\sum_{n\le N}e^{inx}/n$ is bounded. Hence
$F_N(u_x)\ge(4/\pi^2)\log N+O_{x,k,r}(1)$, and
\eqref{eq:cue-leakage} proves that the odd mode tends to zero.
Dominated convergence integrates over $x$. Bounds for the smooth
factor are uniform on every indicated range of $r$.

For the joint assertion fix $s>0$ and let $J$ be the arc of length
$2\pi s/N$ following $x$. In the determinant replace $u_x$ by
$w=u_x z^{\mathbf1_J}$, $|z|\le1$. Then
\[
 \|(w-u_x)Q_{N,x}\|_{\rm HS}^2
 \le4\operatorname{Tr}(Q_{N,x}\mathbf1_JQ_{N,x})\le4s.
\]
For a contraction multiplier $w$, the operator
$I-(QwQ)^*(QwQ)$ is the sum of $Q(1-|w|^2)Q$ and
$Qw^*(1-Q)wQ$. Its determinant therefore has the same bound as
\eqref{eq:cue-leakage} with exponent at least
\[
 \bigl(\|(1-Q)u_xQ\|_{\rm HS}-2\sqrt s\bigr)_+^2.
\]
It tends to infinity uniformly for $z$ on any circle $|z|=b<1$.
Cauchy coefficient extraction proves vanishing of every odd mixed
mode against $\mathbf1_{\#J=l}$, for fixed $l$. The event
$D_{j,N}>s$ is the union $\#J\le j-1$, so all its nonconstant
mixed modes vanish. The even modes already vanish by rotation.

Rescaling by $N/(2\pi)$, and removing a scalar phase, makes
$Q_{N,x}$ converge on compact intervals to $K^0$. Its count
generating functions converge: in the Fredholm series Hadamard's
bound gives a summable majorant $s^m/m!$. The limiting gap has no
atoms, since the intensity is continuous, and the finite gaps are
tight because cyclic summation gives $\mathbb ED_{j,N}=j$.
These facts and the mixed Fourier identities identify the joint
limit as the product measure in the theorem.

Finally,
\[
 \mathbb E\sup_x|\phi'_{N,r}(x)|
 \le\sum_{m\ge1}r^m\mathbb E|\operatorname{Tr}U_N^m|
 \le\sum_{m\ge1}r^m\sqrt m<\infty.
\]
The identity $\mathbb E|\operatorname{Tr}U_N^m|^2=\min(m,N)$ follows
by integrating the two-point determinant against the Fourier factor
$e^{im(\theta-\psi)}$.
Thus the phase change over a tight scaled gap tends to zero in
probability, and
\[
 W_{N,I+j}-W_{N,I}=j-D_{j,N}+o_{\mathbb P}(1).
\]
For uniform $V$ modulo two, the signs at $V$ and $V+v$ agree with
probability $1-d_2(v)$. Their discontinuity sets have zero limiting
measure, which completes the proof.
\end{proof}

\subsection{Distant labels}
A local central limit estimate supplies the fixed-frequency information
needed when the label separation grows.

\begin{theorem}\label{thm:palm-local-limit}
Let $f_j$ be the density of $D_j$, extended by zero on the negative
half-line, and let $\sigma_j=\sqrt{\log j}/\pi$, $j\ge2$. Then
\[
 \int_{\mathbb R}\left|\sigma_j f_j(j+\sigma_j y)
       -\frac{e^{-y^2/2}}{\sqrt{2\pi}}\right|dy\longrightarrow0.
\]
For every fixed $b\ne0$, $\mathbb Ee^{ib(D_j-j)}\to0$.
In particular the correlations in Theorem~\ref{thm:cue-labels}
tend to zero as $j\to\infty$, after taking its $N\to\infty$ limit.
\end{theorem}

\begin{proof}
The sine projection compressed to $(0,s)$ is a positive trace-class
contraction $A_s$. Its count generating function
$\det(I+(z-1)A_s)$ identifies the count as a sum of independent
Bernoulli variables with parameters the eigenvalues of $A_s$.
Direct integration gives
\[
 \operatorname{Tr}A_s=s,\quad
 \operatorname{Tr}(A_s-A_s^2)
 =s-2\int_0^s(s-u)k(u)^2du=\pi^{-2}\log s+O(1).
\]
Indeed $\int_0^\infty k(u)^2du=1/2$; the tail multiplied by $s$
is bounded, and $2\int_0^s uk(u)^2du=\pi^{-2}\log s+O(1)$.
Conditioning at zero and at $s$ removes two evaluation vectors.
The compressed two-point Palm operator $B_s$ therefore satisfies
$R=A_s-B_s\ge0$, $\operatorname{Tr}R\le2$, and
\[
 \operatorname{Tr}B_s=s+O(1),\qquad
 \operatorname{Tr}(B_s-B_s^2)=\pi^{-2}\log s+O(1).
\]
The variance difference is $\operatorname{Tr}R-
\operatorname{Tr}(R(A_s+B_s))$, of modulus at most two.

For completeness, a Bernoulli sum $X$ of mean $\mu$ and variance $v$
satisfies, uniformly over its parameters as $v\to\infty$,
\begin{equation}\label{eq:bernoulli-local}
 \sup_{l\in\mathbb Z}\left|\sqrt v\,\mathbb P(X=l)
  -(2\pi)^{-1/2}e^{-(l-\mu)^2/(2v)}\right|\longrightarrow0.
\end{equation}
Its centered characteristic function obeys
$|\psi(\theta)|\le e^{-2v\sin^2(\theta/2)}
\le e^{-2v\theta^2/\pi^2}$ for $|\theta|\le\pi$.
Taylor expansion of its logarithm gives
$\log\psi(u/\sqrt v)=-u^2/2+O(|u|^3/\sqrt v)$; the cubic
remainder is bounded by a constant times each Bernoulli variance.
Fourier inversion and dominated convergence, after extending the
scaled integrand by zero outside $[-\pi\sqrt v,\pi\sqrt v]$,
prove \eqref{eq:bernoulli-local}. Finite truncation proves the same
assertion for countable Bernoulli sums of finite mean.

The Palm Campbell formula gives
\[
 f_j(s)=(1-k(s)^2)\,
      \mathbb P^{0,s}\{N(0,s)=j-1\},\qquad s>0.
\]
It follows by counting a point at $s$ exactly when there are $j-1$
points between zero and $s$. Applying \eqref{eq:bernoulli-local} to
$B_s$ proves the stated density convergence uniformly on every
bounded interval of $y$. The one-point Palm count has the same
mean and variance estimates. Chebyshev at $s=j\pm M\sigma_j$
gives
\[
 \limsup_{j\to\infty}\mathbb P(|D_j-j|>M\sigma_j)\le C/M^2.
\]
There are infinitely many positive points almost surely: Chebyshev
at dyadic lengths and Borel--Cantelli prove this. Thus $f_j$ has
mass one, and the tail estimate upgrades compact convergence to
$L^1$ convergence. If $g_j(y)=\sigma_jf_j(j+\sigma_jy)$ and
$\varphi$ is the standard Gaussian density, then
\[
 \left|\mathbb Ee^{ib(D_j-j)}-e^{-b^2\sigma_j^2/2}\right|
 \le\|g_j-\varphi\|_1\longrightarrow0.
\]
Taking $b=\pi l$, $l\in\mathbb Z\setminus\{0\}$, proves Haar
convergence modulo two. The continuous function $1-2d_2$ has Haar
mean zero. Since both limiting signs are fair and their joint law
is invariant under flipping both signs, zero correlation makes
all four joint probabilities tend to $1/4$.
\end{proof}

\subsection{The adjacent probability and its certificate}
Let $K_s$ be the sine projection compressed to an interval of length
$s$, and let $E(s)=\det(I-K_s)$ be its empty-interval probability.

\begin{theorem}\label{thm:cue-agreement}
The adjacent agreement probability in Theorem~\ref{thm:cue-labels} is
\[
 A_{\rm adj}=1-2\sum_{n\ge1}(-1)^{n+1}E(n).
\]
The series converges absolutely, and rigorous bounds are
\[
 0.6665466018372107200719020708703
 <A_{\rm adj}<
 0.6665466018372107200719020748216.
\]
\end{theorem}

\begin{proof}
Differentiating the right endpoint of an empty interval and using
Campbell's formula gives $-E'(s)=\mathbb P^0(D_1>s)$, and hence
$E''=f_1$, distributionally. The probability of two points in an
added interval of length $h$ is $O(h^2)$, since the two-point
intensity is bounded by one; this justifies the first derivative.
At zero, $E(0)=1$ and $E'(0)=-1$. The eigenvalues of $K_s$ lie in
$[0,1]$ and sum to $s$, so $0\le E(s)\le e^{-s}$.
Put $g(s)=d_2(1-s)$. Then $g(0)=1$, $g'(0+)=-1$, and
$g''=2\sum_{n\ge1}(-1)^{n+1}\delta_n$ on $(0,\infty)$.
Twice integrating by parts gives
\[
 \mathbb E d_2(1-D_1)=\int_0^\infty g(s)E''(s)ds
      =2\sum_{n\ge1}(-1)^{n+1}E(n).
\]
The boundary terms at zero cancel. Those at infinity vanish because
$E$ decays exponentially and $E'(s)\to0$. Absolute convergence and
monotonicity of $E$ give the alternating upper and lower bounds.

Here is the analytic error estimate used for the numerical interval.
On $a,\xi\in[-1/2,1/2]$ the Fourier-feature kernel is
$T(a,\xi)=\sqrt s\,e^{2\pi is a\xi}$ and $K_s$ is unitarily
equivalent to $TT^*$. Let $T_p$ replace the exponential by its
Taylor polynomial of degree $p$. With
\[
 h=\pi s/2,\quad R=e^h h^{p+1}/(p+1)!,\quad
 \delta=sR(2+R),
\]
we have
\[
 \|TT^*-T_pT_p^*\|_1\le\delta,
 \qquad \operatorname{Tr}(T_pT_p^*)\le s(1+R)^2.
\]
For positive finite matrices, differentiating the determinant along
their positive interpolation, and using the adjugate bound
$\|\operatorname{adj}(I-C)\|\le\prod_l(1+\lambda_l(C))
\le e^{\operatorname{Tr}C}$, proves
\[
 |\det(I-A)-\det(I-B)|\le\|A-B\|_1
           e^{\max(\operatorname{Tr}A,\operatorname{Tr}B)}.
\]
Finite-rank approximation and the absolutely convergent
exterior-power series extend this to positive trace-class operators.
Thus the more generous error
$\delta\exp(1+s+s(1+R)^2)$ encloses the determinant error.

The finite calculation uses orthonormal Legendre polynomials
$L_n(u)=\sqrt{2n+1}\,P_n(2u)$ on $[-1/2,1/2]$. Their exact
moments, from Rodrigues' formula and integration by parts, are
\[
 M_{kn}:=\int_{-1/2}^{1/2}u^kL_n(u)du
 =\frac{\sqrt{2n+1}\,k!}
        {2^k(k-n)!!(k+n+1)!!}
\]
when $k\ge n$ and $k-n$ is even, and zero otherwise.
The matrix of the truncated feature map, divided by $\sqrt s$, is
\[
 T^{(p)}_{mn}=\sum_{k=0}^p
                 \frac{(2\pi is)^k}{k!}M_{km}M_{kn}.
\]
Its even block is real and its odd block purely imaginary, so its
Gram determinant is the product of two real determinants. This
provides a finite expression using rational numbers, square roots
and $\pi$, with no quadrature. Ball arithmetic at 180 decimal digits,
with $p=110$, evaluates $E(1),\ldots,E(7)$ and enlarges every ball
by the analytic error just proved. The alternating sums through
seven and six give, respectively, the asserted lower and upper
bounds. The program \texttt{certify\_cue\_agreement}, its complete
ball intervals and the exact command are available through the
workbench.
\end{proof}
